% -----------------------------------------------------------------------------
% This file is automatically generated.
% Do not edit manually.
% Source: notebooks/support/regression-analysis/support.Rmd
% -----------------------------------------------------------------------------


\noindent We compare the performance of models with four different statistics, the Multiple $r^2$, Adjusted $r^2$ (both obtained from the \ttblue{summary} function), and the AIC and BIC values. These last two can be called with their own respective commands.

\par\addvspace{\topsep}
\begingroup
\fvset{listparameters={%
  \setlength{\topsep}{0pt}%
  \setlength{\partopsep}{0pt}%
  \setlength{\parsep}{0pt}%
  \setlength{\itemsep}{0pt}%
}}
\begin{Verbatim}[breaklines=true,formatcom=\color{ada_blue}]
AIC(ussteamco.mod.0)
\end{Verbatim}
\begin{Verbatim}[breaklines=true,formatcom=\color{ada_light_blue}]
[1] 1202.633
\end{Verbatim}
\begin{Verbatim}[breaklines=true,formatcom=\color{ada_blue}]
AIC(model_backward)
\end{Verbatim}
\begin{Verbatim}[breaklines=true,formatcom=\color{ada_light_blue}]
[1] 1150.277
\end{Verbatim}
\begin{Verbatim}[breaklines=true,formatcom=\color{ada_blue}]
AIC(ussteamco.mod.1)
\end{Verbatim}
\begin{Verbatim}[breaklines=true,formatcom=\color{ada_light_blue}]
[1] 1151.848
\end{Verbatim}
\begin{Verbatim}[breaklines=true,formatcom=\color{ada_blue}]
AIC(ussteamco.mod.2)
\end{Verbatim}
\begin{Verbatim}[breaklines=true,formatcom=\color{ada_light_blue}]
[1] 1143.217
\end{Verbatim}
\endgroup
\par\addvspace{\topsep}
\noindent The AIC values shown above can be compared directly because all four models are fitted to the same response data (\ttblue{USSteamCoEstim}). Model \ttblue{mod.3} is fitted after winsorizing observation 22 and therefore uses a modified response variable. Since the likelihood is computed from the observed responses, the AIC (and BIC) of \ttblue{mod.3} is not directly comparable with those of the earlier models. Instead, \ttblue{mod.3} should be regarded as a sensitivity analysis demonstrating the effect of mitigating the influence of an outlying observation.

\noindent The AIC value of the \ttblue{model\\_backward} does not quite match the output of the \ttblue{step} function. This can be explained as follows.

\noindent The Akaike Information Criterion (AIC) provides a penalized measure of model fit, defined as
$\text{AIC} = -2 \log \hat{L} + 2k$,
where $\log \hat{L}$ is the maximized log-likelihood of the model and $k$ is the number of estimated parameters.

\noindent For a linear regression model with normally distributed errors, this log-likelihood takes the form
$\log \hat{L} = -\frac{n}{2} \left[ \log(2\pi) + 1 + \log\left(\frac{\text{SSE}}{n}\right) \right]$,
leading to the AIC expression
$\text{AIC} = n \left[\log(2\pi) + 1 + \log\left(\frac{\text{SSE}}{n}\right)\right] + 2k$.
In contrast, the \ttblue{step()} function in R reports AIC values based on a deviance-like approximation, that omits constant terms unrelated to model comparison:
$\text{AIC}_{\text{step}} = n \log\left(\frac{\text{SSE}}{n}\right) + 2k$.
The difference between the two expressions is a constant shift of
$n \cdot (\log(2\pi) + 1) + 2$.
The term $n \cdot (\log(2\pi) + 1)$ reflects the part of the log-likelihood that depends only on the sample size and the assumption of normally distributed errors.
The additional +2 arises because R’s AIC() function evaluates the likelihood based on a model with an estimated variance parameter $\sigma^2$, introducing one more estimated parameter than is typically counted in the model formula (e.g., intercept and slopes). This shifts the effective count of parameters k by 1 and contributes an extra penalty of 2 in the AIC expression.

\noindent Similarly, the BIC values are obtained as follows.
